% !TEX root = ../main.tex
\chapter{METODOLOGI PENELITIAN}
\label{bab:metodologi}
% Pedoman 9.3: cukup rinci agar eksperimen dapat diulang oleh orang lain.

\section{Jenis dan Desain Penelitian}
\label{sec:desain}
% Kuantitatif, eksperimen komparatif. Sertakan gambar alur penelitian.
...

\section{Rancangan Sistem}
\label{sec:rancangan}
\subsection{Arsitektur Sistem}
% Alur n8n: Direktur -> Pencari -> Peneliti -> Agen Verifikasi -> Pengamat Teknis.
...
\subsection{Lapisan \asing{Serving}}
% vLLM, model dan kuantisasi, cara tiga model berbagi satu GPU.
...
\subsection{Konfigurasi Agen Verifikasi}
\begin{table}[htbp]
  \caption{Konfigurasi Agen Verifikasi yang dibandingkan}
  \label{tab:konfigurasi}
  \centering\setstretch{1}
  \begin{tabularx}{\textwidth}{@{}l X X@{}}
    \toprule
    Kode & Konfigurasi & Fungsi \\
    \midrule
    S1--S3 & Setiap model kandidat sendirian & \asing{Baseline}; mencari \asing{single model} terbaik \\
    SC & \asing{Single model} terbaik, $k$ sampel, \asing{majority vote} & Kontrol: komputasi setara \asing{council} tanpa keberagaman model \\
    C-vote & Tiga model berbeda, \asing{majority vote} & \asing{Council} paling sederhana \\
    C-chair & Tiga model menjawab, satu \asing{chairman} menyintesis & Pola \asing{LLM council} \\
    \bottomrule
  \end{tabularx}
\end{table}
\subsection{\asing{Prompt} dan Format Keluaran}
...

\section{Objek Penelitian dan \asing{Dataset}}
\label{sec:dataset}
...

\section{Lingkungan dan Instrumen Penelitian}
\label{sec:lingkungan}
...

\section{Prosedur Eksperimen}
\label{sec:prosedur}
...

\section{Metode Analisis Data}
\label{sec:analisis}
% Contoh rumus:
\begin{equation}
  F_1 = \frac{2 \cdot \mathit{precision} \cdot \mathit{recall}}{\mathit{precision} + \mathit{recall}}
  \label{eq:f1}
\end{equation}
...

\section{Etika Penelitian}
\label{sec:etika}
...
